% -----------------------------------------------------------------------------
%  entries/qNNN.tex
%
%  One file per problem, named after its permanent ID.  The ID never changes
%  and never gets reused, even if the entry later moves to a different day.
%
%  Keys are comma-separated, so any value containing a comma must be braced:
%      topic = {Trace norm, Hilbert-Schmidt norm}
% -----------------------------------------------------------------------------
\begin{entry}{Q028}{Three qubits, two ways}

\entrymeta{
  topic   = {Multipartite entanglement, SLOCC classification},
  est     = $\sim$90 minutes,
  origin  = Adapted,          % Sourced / Extracted / Adapted / Original
  source  = {W.~D\"ur, G.~Vidal, J.~I.~Cirac, Phys.\ Rev.\ A \textbf{62}, 062314 (2000)},
  url     = https://arxiv.org/abs/quant-ph/0005115,
  prereq  = Schmidt decomposition,
  status  = drafted,           % drafted / solved / written-up / published
}

\begin{problem}
For two parties sharing a pure state, stochastic LOCC (SLOCC)---local
operations and classical communication that are only required to succeed
with nonzero probability---sorts all states by a single integer, the Schmidt
rank. Three qubits behave differently: there are two inequivalent kinds of
genuinely tripartite entanglement, represented by the GHZ and W states. This
problem proves their inequivalence with a tripartite analogue of the Schmidt
rank, and then shows that the two classes are nevertheless not separated from
each other as subsets of state space.

Work in $\mathcal{H} = (\mathbb{C}^2)^{\otimes 3}$ with computational basis
$\{|jkl\rangle\}_{j,k,l\in\{0,1\}}$, and let
\begin{equation}\label{eq:3q-states}
  |\mathrm{GHZ}\rangle = \tfrac{1}{\sqrt2}\big(|000\rangle + |111\rangle\big),
  \qquad
  |W\rangle = \tfrac{1}{\sqrt3}\big(|001\rangle + |010\rangle + |100\rangle\big).
  \label{q028: GHZ and W states}
\end{equation}
Write $GL(2,\mathbb{C})$ for the group of invertible $2\times 2$ complex
matrices. Two nonzero vectors $|\psi\rangle, |\phi\rangle \in \mathcal{H}$
are \emph{SLOCC-equivalent} if there exist $A, B, C \in GL(2,\mathbb{C})$ and
$\lambda \in \mathbb{C}\setminus\{0\}$ such that
\begin{equation}\label{eq:3q-slocc}
  |\psi\rangle = \lambda\,(A\otimes B\otimes C)\,|\phi\rangle .
\end{equation}
D\"ur, Vidal and Cirac show that this coincides with interconvertibility by
LOCC with nonzero success probability; here \eqref{eq:3q-slocc} is taken as
the definition. The \emph{SLOCC class} of $|\phi\rangle$ is the set of unit
vectors SLOCC-equivalent to it.

The \emph{tensor rank} of $|\psi\rangle \in \mathcal{H}$ is
\begin{equation}\label{eq:3q-rank}
  \operatorname{rk}(\psi)
  = \min\Big\{ r \in \mathbb{N} \;:\;
    |\psi\rangle = \sum_{i=1}^{r} a_i \otimes b_i \otimes c_i,\;
    a_i, b_i, c_i \in \mathbb{C}^2 \Big\},
\end{equation}
with $\operatorname{rk}(0) = 0$. The vectors in a decomposition need not be
normalized or orthogonal.

You may use the following fact. A vector
$v = \sum_{j,k} M_{jk}|jk\rangle \in \mathbb{C}^2\otimes\mathbb{C}^2$ can be
written as $b\otimes c$ (zero allowed) if and only if its coefficient matrix
has rank at most one, that is,
\begin{equation}\label{eq:3q-det}
  v = b\otimes c \text{ for some } b, c
  \quad\Longleftrightarrow\quad
  \det M = 0 .
  \label{q028: identity for rank 2}
\end{equation}
This holds because the rank of $M$ equals the Schmidt rank of $v$ (singular
value decomposition; Nielsen and Chuang, Thm.~2.7).

\begin{enumerate}
  \item[(a)] Show that $\operatorname{rk}$ is invariant under SLOCC
  equivalence \eqref{eq:3q-slocc}, and that
  $\operatorname{rk}(\mathrm{GHZ}) = 2$.

  \item[(b)] Show that $\operatorname{rk}(\mathrm{W}) = 3$

  \item[(d)] Show that $|W\rangle$ is nonetheless a limit of unit vectors in
  the SLOCC class of $|\mathrm{GHZ}\rangle$. Proceed as follows:
  Let $w : \mathbb{R} \to \mathbb{C}^2$, $t \mapsto w(t)$, be a
  differentiable curve of (not necessarily normalized) qubit vectors, with
  the derivative $w'(t)$ taken component-wise. For every $\varepsilon \neq 0$
  such that $w(\varepsilon)$ and $w(0)$ are linearly independent, define
  \begin{equation}\label{eq:3q-quotient}
    |\chi_\varepsilon\rangle
    = \frac{w(\varepsilon)^{\otimes 3} - w(0)^{\otimes 3}}{\varepsilon}.
  \end{equation}
  \begin{enumerate}
    \item[(i)] Show that $|\chi_\varepsilon\rangle \neq 0$ and that
    $|\chi_\varepsilon\rangle / \||\chi_\varepsilon\rangle\|$ lies in the
    SLOCC class of $|\mathrm{GHZ}\rangle$.
    \item[(ii)] Compute $\lim_{\varepsilon\to 0} |\chi_\varepsilon\rangle$ in
    terms of $w(0)$ and $w'(0)$.
    \item[(iii)] Choose $w$ so that this limit is a nonzero multiple of
    $|W\rangle$ and so that $w(\varepsilon)$ and $w(0)$ are linearly
    independent for all sufficiently small $\varepsilon \neq 0$. Conclude.
  \end{enumerate}

\end{enumerate}
\end{problem}

\begin{solution}[a]
Take a vector $v$ whose tensor rank is $r$, then for arbitrary $A, B, C \in GL(2,\mathbb{C})$ and
$\lambda \in \mathbb{C}\setminus\{0\}$ defined the transformed vector $v' = \lambda(A\otimes B\otimes C)v$.
Since $v$ is of rank $r$, we can write it as 
\begin{equation}
    v= \sum_{i=1}^ra_i\otimes b_i\otimes c_i
\end{equation}
thus
\begin{equation}
    v' =\lambda\sum_{i=1}^{r} Aa_i\otimes Bb_i\otimes Cc_i.
\end{equation}
The expression directly shows that the transformation cannot increase the rank, as it is already written as a sum of $r$ terms, so $\operatorname{rk}(v')\leq r$.

Meanwhile, since $A,B,C$ are invertible, we can take $A^{-1},B^{-1},C^{-1} \in GL(2,\CC)$ and $\lambda^{-1} \in \CC \setminus \{0\}$, which itself defines a transformation in the same SLOCC class such that $v = \lambda^{-1}(A^{-1}\otimes B^{-1}\otimes C^{-1})v'$. Analogously, this implies $\operatorname{rk}(v)=r\leq \operatorname{rk}(v')$. Therefore, $\operatorname{rk}(v')=r$ also, being the tensor rank invariant by such transformations.

For the $\ket{GHZ}$ state, it is already written as the sum of two tensors, thus $\operatorname{rk}(\ket{GHZ})\leq 2$

Suppose by contradiction that $\operatorname{rk}(\ket{GHZ})= 1$, thus, there is $a,b,c\in \CC^{2}$ such that $\ket{GHZ}=\ket{a}\otimes\ket b\otimes\ket c$. Define $\ket{a^{\perp}}$ as the normalized vector orthogonal to $\ket{a}$ (which exists and, since the subspace has dimension 2, it is unique), hence
\begin{equation}
    \bra{a^{\perp}}\otimes\id \otimes \id\ket{GHZ} = 0
\end{equation}
However, from the expression \eqref{q028: GHZ and W states},
\begin{equation}
    \bra{a^{\perp}}\otimes\id \otimes \id\ket{GHZ} = \frac{1}{\sqrt{2}} \left( \braket{a^{\perp}}{0} \ket{00} + \braket{a^{\perp}}{1} \ket{11}\right)
\end{equation}
which can never be zero for any non-null vector $\ket{a^{\perp}}$, since the subspace has dimension 2; there is no other direction that is mutually orthogonal to $\ket{0}$ and $\ket 1$. Therefore, by contradiction, we have that $\operatorname{rk}(\ket{GHZ})=2.$

\end{solution}

\begin{solution}[b]
If we assume that $\operatorname{rk}(\ket{W})=1$, we arrive at a similar contradiction as in the previous case. 
We thus assume, again by contradiction, that $\operatorname{rk}(\ket{W})=2$. Hence there are vectors satisfying
$\ket{W}= \ket{a_0,b_0,c_0}+ \ket{a_1,b_1,c_1}$. We must have that at least one pair of vectors from the same subspace is linearly independent (otherwise, if all were linearly dependent, we could reduce it to a separable state $\lambda\ket{a_0,b_0,c_0}$). Without loss of generality, take $\ket{a_0},\ket{a_1}$ to be linearly independent.

Similarly as before we take $\ket{a_i^{\perp}}$ the normalized vector orthogonal to $\ket{a_i}$, such that 
\begin{equation}
     \bra{a_i^{\perp}}\otimes\id \otimes \id\ket{GHZ} = \ket{b_i,c_i}.
\end{equation}

Nonetheless, using the expression from \eqref{q028: GHZ and W states}
\begin{equation}
\begin{split}
    \bra{a_i^{\perp}}\otimes\id \otimes \id\ket{GHZ} &= \frac{1}{\sqrt{3}} (\braket{a_i^{\perp}}{0} \ket{01} + \braket{a_i^{\perp}}{0} \ket{10} + \braket{a_i^{\perp}}{1} \ket{00})\\
    &= \sum_{ij} M^{(i)}_{jk}\ket{jk}
\end{split}
\end{equation}
for the matrix
\begin{equation}
    M^{(i)}=\frac{1}{\sqrt{3}}\begin{pmatrix}
        \braket{a_i^{\perp}}{1} &  \braket{a_i^{\perp}}{0}\\
        \braket{a_i^{\perp}}{0} & 0
    \end{pmatrix}
\end{equation}
whose determinant vanishes if and only if $\ket{a_i^{\perp}}=\ket{1}$, but that would mean $\ket{a_0},\ket{a_1}$ are linearly dependent, which is contradiction. 

Therefore, by exhausting the other possibilities, $\operatorname{rk}(\ket{W})=3$.



\end{solution}




\begin{solution}[c]
Given $\epsilon>0$, since $\ket{0},\ket{1}$ are orthogonal there is a $A_\epsilon \in GL(2,\CC)$ such that $A_\epsilon\ket{0}=w(\epsilon)$ and $A_\epsilon\ket{1}=-w(0)$, for any $w(\epsilon)\in \CC^{2}$ and $w(0)\in \CC^{2}$. Moreover, we can take $\lambda = \sqrt{2}/\epsilon$ such that 
\begin{equation}
    \lambda (A_\epsilon^{\otimes3})\ket{GHZ}= |\chi_\varepsilon\rangle
\end{equation}

The limit can be obtained through the derivative of the function $f(t)=w(t)^{\otimes3}$, which, by the product rule (notice that the proof for tensor product is analogous to regular multiplication of functions), is
\begin{equation}
    f(t) = w(t)\otimes w(t) \otimes w'(t) + w(t)\otimes w'(t) \otimes w(t) + w'(t)\otimes w(t) \otimes w(t)
\end{equation}
such that 
\begin{equation}
    \lim_{\epsilon \rightarrow 0} \ket{\chi_\epsilon}= w(0)\otimes w(0) \otimes w'(0) + w(0)\otimes w'(0) \otimes w(0) + w'(0)\otimes w(0) \otimes w(0)
\end{equation}

This limit is proportional to the $\ket{W}$ state if $w'(0)=\ket{1}$ and $w(0)=\ket{0}$. From the argument that there is always a matrix $A_\epsilon$ satisfying $A_\epsilon\ket{0}=w(\epsilon)$ and $A_\epsilon\ket{1}=-w(0)$, it directly follows that the desired limit exists and thus that the SLOCC class of $\ket {GHZ}$ is not closed.
 
 But we can construct an example explicitly. A natural choice of curve that satisfies these initial conditions is $w(t)=\cos(t)\ket{0}+\sin(t)\ket{1}$. We then need that
 \begin{equation}
     A_\epsilon \ket{0} = \cos(\epsilon) \ket{0} + \sin(\epsilon)\ket{1}, \quad A_\epsilon \ket{1} = -\ket{0}
 \end{equation}
 which is satisfied by the matrix 
 \begin{equation}
     A_\epsilon = \begin{pmatrix}
         \cos(\epsilon)  &  -1 \\
          \sin(\epsilon) & 0
     \end{pmatrix}
 \end{equation}

\end{solution}



% \begin{remark}
% \end{remark}

\begin{takeaway}
\end{takeaway}

%  Everything below is optional.  An entry with no attempt block is complete.
%  Record only what would change what you do next time: typos and arithmetic
%  slips you caught yourself are artefacts of the medium, not facts about the
%  mathematics.  Add "public," to show this block in the public build.
% \attempt{
%   date       = ,               % ISO, YYYY-MM-DD -- required if the block is used
%   time       = ,
%   difficulty = ,               % 1-5, perceived
%   howitwent  = ,
%   insight    = ,
%   note       = ,
% }

\end{entry}